\section{‘return’ and ‘yield’ Statements}
The \lstinline|return|\index{return} statement has to be the last statement of a method, while the \lstinline|yield|\index{switch!yield} statement has to be the last statement of a \lstinline|case|\index{switch!case} block of a \lstinline|switch| expression\index{switch!expression}; refer to the chapters \tqfullvref{sec:ReturningValues} and \tqfullvref{sec:CaseResults} for more details on this.

When \lstinline|return|\index{return} is always the last statement of a method, it means that the simple form
\begin{lstlisting}
return;
\end{lstlisting}
is obsolete and will be omitted. This implies that there is only one \lstinline|return|\index{return} statement per method allowed, and that short-cutting a method is a forbidden practice.

Usually, the name of the return value is \lstinline|retValue| (see \tqfullvref{lbl:RetValue}), and the lines before the \lstinline|return|\index{return} statement are a blank line and the comment \lstinline|//---* Done *---…|:
\keeplisting{7}
\begin{lstlisting}
public final void myMethod()
{
    …
        
    //---* Done *----------------------------------------------------
    return retValue;
}   //  myMethod()
\end{lstlisting}

\lstinline|yield|\index{switch!yield} as a keyword was introduced with the new \lstinline|switch|\index{switch} format and \lstinline|switch| expressions\index{switch!expression} by Java~12. It is used to return a value from a \lstinline|case|\index{switch!case} block that is not an expression. Same as \lstinline|return|\index{return} has \lstinline|yield|\index{switch!yield}to be placed on a line of its own, a the end of a \lstinline|case| block\index{switch!case}. Without  a return value \lstinline|yield|\index{switch!yield} will always cause a compiler error. Some more details can be found at \autocite{ORACLE_DOC_LANGUAGE_SPECIFICATION:Yield}.

The use of round brackets with the \lstinline|return|\index{return} or \lstinline|yield|\index{switch!yield} statement is not required unless they make the return value more obvious and better readable in some way.

\subsubsection{Principles}
\begin{enumerate}[label=P\arabic*.]
\item{The \lstinline|return|\index{return} statement is, if it exists at all, always the last statement in a method.

There is at maximum one \lstinline|return|\index{return} statement in a message.}
\item{The \lstinline|return|\index{return} statement is introduced by a blank line and a \verb#Done!# comment.}
\item{A \lstinline|yield|\index{switch!yield} is always the last statement in a \lstinline|case|\index{switch!case} block.}
\item{The \lstinline|return|\index{return} and \lstinline|yield|\index{switch!yield} statements are always placed on a line of their own.}
\item{Do not use round brackets with \lstinline|return|\index{return} or \lstinline|yield|\index{switch!yield}.}
\end{enumerate}

%==============================================================================
% End of file!!
%==============================================================================
